\begin{definition}[Actual reunion measurement region]%
    \label{def:peps_torus_swept_route_reunion_region}
    \lean{TNLean.PEPS.torusSweptRouteReunionRegion}
    \leanok
    \uses{def:peps_torus_swept_route_partner_walks}
    On a torus with periods at least eight and seven, take the translated
    six-site two-plaquette block whose lower-left site is $(0,-2)$
    relative to the translation vertex. It contains the actual final
    B reunion connector and both partner loops.
\end{definition}

\begin{theorem}[Original-spin measurement of the derived reunion class]%
    \label{thm:peps_torus_swept_route_joint_flux_measurement}
    \lean{TNLean.PEPS.regularWalkHolonomy_torusSweptRoute_jointClass_trivial_iff,
        TNLean.PEPS.IsGIsometric.exists_torusSweptRouteReunion_physicalMeasurement}
    \leanok
    \uses{def:peps_torus_swept_route_reunion_region,
        thm:peps_torus_swept_route_reunion_holonomy,
        thm:peps_regular_inserted_walk_holonomy,
        thm:peps_closed_walk_class_physical_measurement,
        thm:peps_physical_cut_column_action}
    Let every original site tensor be regular G-isometric.
    There is one complete orthogonal projective measurement on the six
    original spins, chosen before both flux labels and all crossing,
    exterior and boundary data, whose conjugacy-class outcome on every
    actual final-route boundary column is the class of
    $ghg^{-1}h^{-1}$. The same eigenaction holds on the actual physical
    cut coefficient matrix. The identity-class outcome occurs precisely
    when $g$ and $h$ commute.
    The additional outcome is the orthogonal complement of the physical
    block image and annihilates these columns.
    This is an auxiliary finite-torus measurement for
    \cite[\S6.6, local lines 2387--2415]{Schuch2010PEPS}; implementation of
    the preceding twelve physical steps is a separate statement.
\end{theorem}

\begin{proof}\leanok
    Restrict the actual reunion walk to its six-site induced graph and
    apply the closed-walk class measurement there. Arbitrary crossing
    and exterior operators do not alter its holonomy. The derived joint
    holonomy is conjugate to the displayed commutator. The all-boundary
    column action gives the physical cut action by finite contraction.
\end{proof}
